\section{Background}
\label{sec:background}

\subsection{Model Context Protocol}

Model Context Protocol (MCP)~\citep{mcp2024} is a JSON-RPC~2.0 protocol that standardizes
how AI agents discover and invoke tools hosted on external servers. A \emph{tool} is a named
function with a JSON Schema describing its parameters. MCP defines three primitive operations:
\texttt{tools/list} to enumerate available tools, \texttt{tools/call} to invoke one, and
\texttt{resources/read} to retrieve structured data.

An MCP session begins when a client establishes a transport connection (Streamable HTTP,
which replaced HTTP with Server-Sent Events in the 2025-03-26 revision, or stdio) to an MCP
server and exchanges \texttt{initialize} messages.
The server advertises its tool catalog; the client may cache it. Subsequent tool calls are
stateless from the protocol's perspective: each \texttt{tools/call} carries the tool name
and a JSON argument object and receives a JSON content array in response.

MCP's authorization specification, added in the 2025-03-26 revision, uses OAuth 2.1 to
authorize a client's access to a server. It does not define a per-call policy decision over
tool arguments or session history, does not require any audit record of calls made, and
does not provide a mechanism for a client to verify that the tool definition it cached at
session start is still the definition the server will execute. These gaps are the governance problem
cMCP addresses.

\subsection{Trusted Execution Environments}

A Trusted Execution Environment (TEE) is a CPU-enforced isolation boundary that provides
two guarantees relevant to this paper: \emph{confidentiality} (memory within the TEE is
inaccessible to software outside, including the hypervisor and OS) and
\emph{remote attestation} (the TEE can produce a hardware-signed measurement report that
proves, to a remote verifier, which code is running inside).

The measurement report is the key primitive. On AMD SEV-SNP~\citep{amdsevsnp2023}, the
report is signed by the AMD attestation key (VCEK) and includes a 64-byte
\texttt{REPORT\_DATA} field that the guest can set to arbitrary bytes. On Intel
TDX~\citep{inteltdx2023}, the equivalent field is the 64-byte \texttt{REPORTDATA}. cMCP uses this field to bind application-level
context (specifically, the session nonce) to the hardware attestation, as described in
Section~\ref{sec:design}.

In software-only mode, cMCP substitutes a software-generated key and produces a synthetic
attestation report marked with \texttt{platform: software-only}. Software-only TRACE Claims
carry \texttt{attestation\_assurance: none} and must not be used for compliance decisions.
All experiments in this paper run in software-only mode.

\subsection{TRACE Standard}

TRACE (Trust Runtime Attestation and Compliance Evidence)~\citep{trace2025} is an open
attestation standard for agentic AI. A TRACE Claim is a JSON document with two top-level
sections: a canonical \emph{trace} envelope (eat\_profile, issued-at timestamp, subject
identifier, runtime measurement, policy reference, data classification, confirmation key)
and a profile-specific addenda section that carries deployment context.

The cMCP profile of TRACE adds a \emph{gateway addenda} section containing: the session
identifier, gateway version, a monotonically increasing sequence number for claim ordering,
an audit chain summary (root hash, tip hash, chain length), a call summary (total calls,
allowed, denied, faulted, tools invoked, maximum session sensitivity), a catalog hash, and
attestation freshness fields. The claim is signed with an Ed25519 key whose public half
appears in the \texttt{cnf} (confirmation key) field, allowing any verifier to check the
signature without contacting the gateway.

\subsection{Cedar Policy Language}

Cedar~\citep{cedar2023} is a policy language and evaluation engine developed by AWS. Cedar
policies express access control decisions as \texttt{permit} or \texttt{forbid} rules over
a triple of \emph{principal} (who is making the request), \emph{action} (what is being
requested), and \emph{resource} (what is being acted on), with optional \texttt{when}
conditions evaluated against a context object.

cMCP evaluates Cedar policies at tool-call time. The principal is the agent's SPIFFE
identity~\citep{spiffe2018}, the action is the tool name, and the resource is the MCP server
URL. The context includes the tool's compliance domain (e.g., \texttt{phi}, \texttt{pii},
\texttt{confidential}), the current session sensitivity level, and the call's input
parameters. Cedar's structural typing ensures that policy evaluation is decidable in
bounded time, a property important for a gateway that must not add unbounded latency.

\subsection{Threat Model}

cMCP assumes the following attacker capabilities:

\textbf{T1: Policy modification.} An attacker with access to the gateway host can modify
the Cedar policy files after TEE startup, attempting to relax enforcement for subsequent
calls.

\textbf{T2: Tool catalog tampering.} An attacker can modify an MCP server's tool
definitions after the gateway has cached them, replacing a benign implementation with one
that exfiltrates data or behaves differently than the approved specification.

\textbf{T3: Claim replay.} An attacker intercepts a TRACE Claim produced for one session
and presents it as evidence for a different session, a different TEE instance, or a future
point in time.

\textbf{T4: Selective disclosure.} An attacker selectively presents audit entries from a
TRACE Claim, omitting evidence of non-compliant calls while presenting evidence of compliant
ones.

\textbf{T5: Cross-boundary data exfiltration.} An agent that has processed regulated data
(PHI, PII) subsequently calls external tools in a way that the orchestrator cannot detect,
because the PHI was in an intermediate agent state rather than in the tool call parameters.

\textbf{T6: Binary substitution.} In a cross-organizational setting, the server-side TEE
binary is swapped after the gateway has verified the server's attestation report, allowing
the new binary to behave differently without re-attestation.

cMCP does not address denial-of-service against the gateway itself, compromise of the CPU
manufacturer's signing key, or attacks that require the TEE's measurement to be forged. It
also does not address semantic attacks where a malicious tool produces a compliant-looking
response that causes the agent to take harmful downstream actions.
